\section{Scaling Laws 的历史脉络}

\subsection{理论根源：统计学习理论}

Scaling Laws 并非全新概念。在经典统计学习理论中，人们早已研究过``模型性能如何随数据量和模型复杂度变化''的问题。

\begin{figure}[H]
\centering
\includegraphics[width=0.95\textwidth]{slides-images/slide-003.jpg}
\caption{统计学习理论中的 Scaling：泛化界与非参数收敛率\protect\footnotemark}
\end{figure}
\footnotetext{Slides 第4页。}

在有限假设类 $\mathcal{H}$ 中选择模型时，经典泛化界给出：

$$\text{Excess Risk} \leq O\left(\sqrt{\frac{\log |\mathcal{H}|}{m}}\right)$$

其中 $m$ 是样本数。这告诉我们误差随样本量以 $1/\sqrt{m}$ 的速率衰减——这本身就是一种``Scaling Law''。

对于更灵活的非参数模型类（如核密度估计），$\beta$-光滑函数的 $L_2$ 估计误差为：

$$\text{Error} \leq O\left(n^{-\frac{\beta}{2\beta+1}}\right)$$

\begin{knowledgebox}{从理论到经验：一个根本性的跃迁}
理论给出的是\textbf{上界}（upper bounds），而非实际实现的损失值。Scaling Laws 的真正突破在于：从理论预测的收敛\textbf{速率}出发，转向\textbf{经验性地拟合}模型性能与资源之间的精确关系。理论提供的是``应该期待什么形式的衰减''，经验拟合则告诉我们``具体的系数和指数是多少''。
\end{knowledgebox}

\subsection{第一篇 Scaling Law 论文：Bell Labs 1993}

Tatsu 认为，历史上第一篇真正的 Scaling Law 论文可能是 1993 年 NeurIPS 上来自 Bell Labs 的工作。论文的作者包括 Vapnik、Cortes 等经典机器学习理论的先驱。

\begin{figure}[H]
\centering
\includegraphics[width=0.95\textwidth]{slides-images/slide-004.jpg}
\caption{Bell Labs 1993：可能是最早的 Scaling Law 论文\protect\footnotemark}
\end{figure}
\footnotetext{Slides 第5页。}

这篇论文的核心思想惊人地超前：
\begin{itemize}
    \item ``训练大规模分类器非常耗费计算资源，需要在训练前预测模型好坏''
    \item 测试误差可表示为 $\text{Error} = \epsilon_\infty + a \cdot n^{-\alpha}$（不可约误差 + 多项式衰减项）
    \item 通过小模型拟合曲线，可以准确预测大模型的行为
\end{itemize}

\begin{knowledgebox}{历史的回响}
这些想法在 30 年后的今天看来如此自然，但在 1993 年就已经被 Bell Labs 的研究者们清晰地表述出来了。正如许多技术突破一样，Scaling Laws 的种子在很早之前就已播下。
\end{knowledgebox}

\subsection{数据缩放的早期研究}

\begin{figure}[H]
\centering
\includegraphics[width=0.95\textwidth]{slides-images/slide-005.jpg}
\caption{Banko \& Brill (2001) 以及 Hestness et al. (2017) 的数据缩放研究\protect\footnotemark}
\end{figure}
\footnotetext{Slides 第6页。}

\textbf{Banko \& Brill (2001)} 研究了 NLP 系统性能如何随数据量缩放。他们发现性能改善是可预测的，并提出了一个至今仍有争议的问题：\textbf{应该花时间改进算法，还是只需收集更多数据？}

\textbf{Hestness et al. (2017)} 可能是大规模神经网络 Scaling Law 的最早系统研究。他们在机器翻译、语音识别和视觉任务上展示了\textbf{误差率以幂律衰减}的规律，并提出了模型行为的三个阶段：

\begin{importantbox}{模型性能的三个阶段}
\begin{enumerate}
    \item \textbf{随机猜测阶段}：数据太少，模型表现接近随机
    \item \textbf{幂律缩放阶段}：性能随资源投入可预测地改善（这是 Scaling Laws 最有用的区间）
    \item \textbf{渐近阶段}：接近模型类的不可约误差，改善越来越慢
\end{enumerate}
\end{importantbox}

\begin{figure}[H]
\centering
\includegraphics[width=0.95\textwidth]{slides-images/slide-006.jpg}
\caption{Hestness et al. (2017) 的三阶段模型以及跨任务的幂律缩放\protect\footnotemark}
\end{figure}
\footnotetext{Slides 第7页。}

Hestness 等人的论文虽然早于后来的大规模 Scaling Law 热潮，但已经包含了许多关键洞察：
\begin{itemize}
    \item 在随机猜测区域用 Scaling Law 做预测很困难（``涌现能力''的前兆）
    \item 可预测的缩放意味着\textbf{计算扩展}非常重要
    \item 量化等效率技术是值得的——如果缩放是可预测的，就应该愿意用计算来换取精度
\end{itemize}

\begin{warningbox}{``涌现能力''与 Scaling Laws 的关系}
近年来关于``涌现能力''（emergent capabilities）的讨论非常热烈，但实际上 Hestness 在 2017 年就已经指出：在模型处于随机猜测区域时，Scaling Laws 的预测是不可靠的。所谓``涌现''，可能部分是从随机猜测区域进入幂律区域时的``跳跃''。
\end{warningbox}

\subsection{本章小结}

Scaling Laws 有深厚的历史根基：从 VC 维和非参数收敛率的理论，到 1993 年 Bell Labs 的开创性工作，再到 2017 年 Hestness 的多任务验证。核心发现是：模型性能与数据/计算之间存在可预测的幂律关系，这一规律跨领域、跨模型族普遍成立。

%% ========== Part 3: 数据 Scaling Laws ==========

